\section{Spectral characterisation of the Picard splitting}
\label{sec:spectral}

The explicit treatment of the nonlinear source gives the Jacobian of Eq.~\eqref{eq:residual} the form of a matrix splitting, with $\mathbf{A}$ of Eq.~\eqref{eq:coefficients} as the splitting operator and the matrix $\mathbf{G}=\mathbf{A}^{-1}\mathbf{D}$ of Eq.~\eqref{eq:picard_operator} as the iteration matrix of the associated stationary method, which is the Picard iteration of Eq.~\eqref{eq:picard_iteration}. The results collected here are specialisations of classical results on matrix splittings and stationary iterations~\cite{Varga2000MatrixAnalysis,Young1971IterativeSystems} to that splitting. They are collected because the splitting is not chosen by the analyst, as Jacobi and Gauss--Seidel splittings are, but imposed by the discretisation.

Two configurations are distinguished. In the general case $\mathbf{D}$ carries the reconstruction stencil and is not symmetric. In the reference configuration of Section~\ref{subsec:picard_iteration}, where the source derivative is evaluated pointwise, $\mathbf{D}$ is diagonal with entries $-\left.\dv*[2]{F}{\Psi}\right|_{P}\Delta V_{P}$, and is therefore positive definite when $\dv*[2]{F}{\Psi}<0$ at every generator. Everything below is local: it refers to a fixed solution $\bm{\Psi}^{\ast}$ and to the Picard map $\bm{\Phi}$ of Eq.~\eqref{eq:picard_iteration} at the frozen threshold of Section~\ref{subsec:picard_iteration}, which is continuously differentiable near $\bm{\Psi}^{\ast}$ with derivative $\mathbf{G}$ there. Throughout, $\mathbf{A}$ is symmetric positive definite and therefore nonsingular, as established in Section~\ref{subsec:algebraic_system}.

\subsection{Local convergence and regularity}
\label{subsec:convergence_regularity}

By Ostrowski's theorem, $\bm{\Psi}^{\ast}$ is a point of attraction of the Picard iteration if $\rho(\mathbf{G})<1$~\cite[Theorem~10.1.3]{Ortega2000IterativeVariables}, with asymptotic convergence factor $\rho(\mathbf{G})$~\cite{Ortega2000IterativeVariables}, and it is not a point of attraction if $\rho(\mathbf{G})>1$~\cite{Ostrowski1973SolutionSpaces}; at $\rho(\mathbf{G})=1$ the linearisation is inconclusive.

In the reference configuration with $\dv*[2]{F}{\Psi}<0$ at every generator, $\mathbf{A}$ and $\mathbf{D}$ are symmetric positive definite, so $\mathbf{G}$ is similar to $\mathbf{S}=\mathbf{A}^{-1/2}\mathbf{D}\mathbf{A}^{-1/2}$ and $\mathbf{J}=\mathbf{A}^{1/2}(\mathbf{I}-\mathbf{S})\mathbf{A}^{1/2}$ is congruent to $\mathbf{I}-\mathbf{S}$. Hence $\sigma(\mathbf{G})$ is real and positive, $\rho(\mathbf{G})=\lambda_{\max}(\mathbf{S})$, and $\rho(\mathbf{G})<1$ if and only if $\mathbf{J}$ is positive definite (cf.~\cite{Householder1964TheAnalysis,Ortega2000IterativeVariables}); in particular, $\rho(\mathbf{G})$ reaches unity exactly where $\mathbf{J}$ becomes singular. With the reconstructed source, $\mathbf{D}$ is not symmetric and $\sigma(\mathbf{G})$ may contain complex pairs, so $\rho(\mathbf{G})$ can cross unity with $\mathbf{J}$ nonsingular, and only the implication from $\rho(\mathbf{G})<1$ to a nonsingular $\mathbf{J}$ of Section~\ref{subsec:newton_method} remains.

\subsection{Krylov iterations under the invariant preconditioner}
\label{subsec:krylov}

In the reference configuration the preconditioned operator of Eq.~\eqref{eq:preconditioned_newton_spectrum} is self-adjoint and positive definite in the inner product induced by $\mathbf{A}$ whenever $\rho(\mathbf{G})<1$, with condition number $(1-\mu_{\min})/(1-\mu_{\max})$, where $\mu_{\min}$ and $\mu_{\max}$ are the extreme eigenvalues of $\mathbf{G}$, and the standard bound for the preconditioned conjugate gradient method applies~\cite{Saad2003IterativeSystems}.

No corresponding statement is made in the general case. There $\mathbf{A}^{-1}\mathbf{J}$ need not be normal, and for any prescribed spectrum there exist operators with that spectrum and arbitrarily poor GMRES convergence~\cite{Greenbaum1996AnyGMRES}. The clustering of $\{1-\mu\}$ about unity therefore gives no bound on the number of GMRES iterations a Newton correction requires. For the same reason $\rho(\mathbf{G})$ fixes the asymptotic factor of the Picard iteration but not its transient, since $\|\mathbf{G}^{k}\|$ may exceed $\rho(\mathbf{G})^{k}$ for small $k$.